% !PRINTED 395
\weberpage{395}
\begin{center}
\S.\,127.

Die Pell'sche Gleichung.
\end{center}

Ist $\omega$ eine zur Discriminante $D$ gehörige reducirte Zahl, so
erhält man durch die Kettenbruchentwickelung
\weq{1}{\omega = \dfrac{P_n \omega_n + P_{n-1}}{Q_n \omega_n + Q_{n-1}},}
und es ist $\omega_n = \omega$ immer dann, wenn $n$ ein Vielfaches von der
Gliederzahl der Periode $\nu$ ist. Dann aber folgt aus (1):
\weq{2}{Q_n \omega^2 = (P_n - Q_{n-1}) \omega + P_{n-1}.}
Es genüge nun $\omega$ der quadratischen Gleichung \S.\,122, (5):
\weq{3}{c \omega^2 = a + b \omega, \quad D = b^2 + 4 a c,}
und aus dieser muss, da $a$, $b$, $c$ ohne gemeinsamen Theiler
sind, die Gleichung (2) durch Multiplication mit einer ganzen
Zahl folgen. Bezeichnen wir diese ganze Zahl mit $u$, so ist also
\weq{4}{Q_n = u c, \quad P_{n-1} = u a, \quad P_n - Q_{n-1} = u b.}
Setzen wir noch
\weq{5}{P_n + Q_{n-1} = t,}
so dass auch $t$ eine ganze positive Zahl ist, so folgt aus (4) und (5)
\[
P_n = \frac{t + u b}{2}, \quad Q_n = u c,
\]
\weq{6}{P_{n-1} = u a, \quad Q_{n-1} = \frac{t - u b}{2},}
und also mit Rücksicht auf die Gleichung
\[
P_n Q_{n-1} - Q_n P_{n-1} = (-1)^n\text{:}
\]
\weq{7}{t^2 - D u^2 = (-1)^n 4.}
Die Gleichung (7) heisst die \gesperrt{Pell'sche Gleichung}. Die
Aufgabe ist die, für ein gegebenes $D$ alle ihre ganzzahligen
Lösungen $t$, $u$ zu finden. Da, wenn $t$, $u$ der Gleichung genügen,
auch $\pm t$, $\pm u$ ihr genügen, so können wir uns auf die Ermittelung
ihrer positiven Lösungen beschränken.

Die Formeln (4), (5) geben uns eine unendliche Zahl solcher
Lösungen; wir haben nur für $n$ ein beliebiges Vielfaches der
Periodenzahl $\nu$ zu wählen, für $u$ den grössten gemeinschaftlichen
Theiler von $Q_n$, $P_{n-1}$, $P_n - Q_{n-1}$, oder auch einfach die Zahl
%CONT>>
